% ===========================================================================
% REVISIONS TO THE MANUSCRIPT TEXT ("poisson_paper.tex")
% G. Garg and R. Mishra
%
% This file contains drop-in replacements for every passage whose claims
% were contradicted (or left unverifiable) by the corrected theory
% (asimex_proofs.tex) and the completed simulations and application
% (results/section5_simulations.tex, results/section6_application.tex).
% Each block states exactly what it replaces. A change log is at the end.
%
% Editorial conventions: theorem names rather than numbers are used where
% the reference crosses documents ("the consistency theorem", "the
% growing-degree central limit theorem"); fix numbering at merge time.
% ===========================================================================

% ---------------------------------------------------------------------------
% (1) REPLACES THE ABSTRACT (page 1 of the draft)
% ---------------------------------------------------------------------------
\begin{abstract}
Poisson regression is widely used for count outcomes, but covariates are
frequently observed with classical measurement error, and the resulting
naive estimates are attenuated with invalid confidence intervals. We
develop Adaptive SIMEX (A-SIMEX), which automates the two ad-hoc choices
of SIMEX---the extrapolation region and the polynomial degree---by
cross-validation with a one-standard-error parsimony rule, adds a
ridge-corrected second extrapolation stage, and is accompanied by
complete asymptotic theory: consistency under a growing-degree regime
that we make explicit; asymptotic normality with an exact decomposition
of the asymptotic variance into data and simulation components (the
latter vanishing at rate $B^{-1}$ in the number of simulation
replicates); and efficiency results relative to the oracle estimator,
including an information-monotonicity lemma showing that no estimator
based on noisy covariates can improve on the oracle to first order,
while the second-order efficiency comparison is signed neither way. The
theory corrects several claims of earlier SIMEX treatments and resolves
the tension between consistency (which requires the polynomial degree to
grow like $\log n$) and inference through an explicit window condition
on the degree growth rate. In simulation studies with $1{,}000$ Monte
Carlo repetitions per design cell---including misspecified (Laplace)
error, negative-binomial outcomes, error variance estimated from
replicated measures, and corrected-score and regression-calibration
competitors---A-SIMEX removes most of the naive attenuation with valid
inference when the error is moderate, and we identify an important
limitation of automatic tuning: held-out prediction on the observed
scale systematically under-extrapolates when the error is large,
because its population minimizer is the naive pseudo-parameter. We
characterize the resulting bias--variance trade-off, quantify the
efficiency premium of extrapolation against the oracle, and give
practical remedies. The methods are implemented in an open-source R
package with an analytic sandwich variance estimator that we validate
against the parametric bootstrap, and are applied to NHANES 2017--2018,
where the measurement error is directly measurable: self-reported
weight in a model for the count of prescription medications. All
results are reproducible from a documented archive with fixed seeds.
\end{abstract}

% ---------------------------------------------------------------------------
% (2) REPLACES SECTION 1.2 ("This Paper's Contributions") IN FULL
% ---------------------------------------------------------------------------
\subsection{This Paper's Contributions}\label{sec:contrib}

We address these gaps by developing Adaptive SIMEX (A-SIMEX) and by
subjecting it to a level of scrutiny that we argue should be standard
for bias-correction methods. The contributions are:

\begin{enumerate}\itemsep3pt
\item \textbf{Automated hyperparameter selection, with its limits
      quantified.} Extrapolation region and polynomial degree are
      selected by $k$-fold cross-validation of the held-out Poisson
      deviance, with a one-standard-error parsimony rule. This removes
      subjectivity and is fully reproducible. Our simulations,
      however, reveal a structural limitation that we do not paper
      over: the population minimizer of the observed-scale CV criterion
      is the naive pseudo-parameter $\beta(1)$, not $\beta_0$, so CV
      systematically under-extrapolates when the measurement error is
      large (Section~\ref{sec:sims1}). Automatic tuning is a solution
      for small-to-moderate error and a documented failure mode for
      large error, for which we propose deconvolution-aware remedies.
\item \textbf{A bias-corrected two-stage extrapolation and an exact
      variance estimator.} The primary polynomial fit is stabilized by
      a ridge second stage, and inference uses an analytic sandwich
      that decomposes the asymptotic variance exactly into a
      \emph{data} component and a \emph{simulation} component of order
      $B^{-1}$; this decomposition corrects the additive
      ``parameter-plus-extrapolation'' split that appears in earlier
      SIMEX treatments and in a draft of this paper, and it extends
      robustly to overdispersed outcomes and to estimated error
      variance (Sections~\ref{sec:sims-nb} and~\ref{sec:sims-sigma2}).
\item \textbf{Complete and corrected asymptotic theory.} We prove
      consistency and asymptotic normality of the A-SIMEX estimator,
      a growing-degree central limit theorem that reconciles the
      $\log n$ degree growth needed for consistency with centered
      inference through an explicit window condition on the analyticity
      margin, and oracle-efficiency results comprising an
      information-monotonicity lemma (no estimator based on noisy
      covariates beats the oracle to first order) and an explicit
      second-order expansion whose sign is design-dependent. Along the
      way we correct three claims that circulate informally for SIMEX:
      the exactness of the extrapolation target, the form of the
      variance decomposition, and the sign of the second-order
      efficiency loss.
\item \textbf{Evidence at replication standard.} All simulations
      ($1{,}000$ Monte Carlo repetitions per cell) and the application
      are generated by a single seed-fixed pipeline; the paper reports
      what the pipeline produces, including results that are
      unflattering to the method.
\end{enumerate}

% ---------------------------------------------------------------------------
% (3) REPLACES SECTION 1.3 ("Roadmap") IN FULL
% ---------------------------------------------------------------------------
\subsection{Roadmap}\label{sec:roadmap}

Section~\ref{sec:background} reviews the model and standard SIMEX.
Section~\ref{sec:method} presents A-SIMEX: cross-validated selection of
$(\lambda_{\max},\text{degree})$ (Section~3.1), the ridge second stage
(Section~3.2), and the analytic sandwich variance estimator
(Section~3.3). Section~\ref{sec:theory} states and proves the main
results: the population SIMEX curve and its analyticity (Lemmas on the
pseudo-true parameter), the consistency theorem, the asymptotic
normality theorem with the exact data/simulation variance
decomposition, the growing-degree central limit theorem, and the
oracle-efficiency theorem with the information-monotonicity lemma.
Section~\ref{sec:sims} reports the simulation studies
($1{,}000$ repetitions per cell): the low-dimensional strong-signal
scenario with $n$- and $\sigma^2_\epsilon$-grids, misspecified Laplace
error, computation times up to $p=50$, negative-binomial outcomes, and
estimated error variance from replicated measures. Section~\ref{sec:app}
applies the methods to NHANES 2017--2018, where self-reported weight is
the error-prone covariate and the error variance is measurable within
the same survey, with a sensitivity analysis over
$\sigma^2_\epsilon$. Section~\ref{sec:software} describes the R
package. Section~\ref{sec:discussion} discusses limitations and
extensions. Complete proofs are collected in a companion appendix
(reproduced in the replication archive as
\texttt{asimex\_proofs.pdf}), which supersedes the proof sketches of
the earlier draft.

% ---------------------------------------------------------------------------
% (4) REPLACES SECTION 3.3 ("Variance Estimation and Confidence
%     Intervals") IN FULL.  The draft's eqs. (14)-(19) and the
%     decomposition Var_param + Var_extrap + Var_sim are superseded.
% ---------------------------------------------------------------------------
\subsection{Variance Estimation and Confidence Intervals}\label{sec:var}

Let $c=(c_0,\dots,c_{K-1})$ denote the combined extrapolation weights
(primary polynomial plus ridge second stage), so that the A-SIMEX
estimator is exactly the linear functional
$\hat\beta_{\asimex}=\sum_jc_j\bar\beta(\lambda_j)$ with
$\sum_jc_j=1$, and let
$\bm H_j=\E\bigl[e^{\beta(\lambda_j)^\top W(\lambda_j)}
W(\lambda_j)W(\lambda_j)^\top\bigr]$ be the per-observation Hessian of
the population SIMEX objective at grid point $\lambda_j$. Under the
conditions of the asymptotic normality theorem,
\begin{equation}\label{eq:sandwich}
  \Var\bigl(\hat\beta_{\asimex}\bigr)\;\approx\;\frac1n\,\bm C
  \bigl(\btV^{\mathrm{data}}+\tfrac1B\,\btV^{\mathrm{sim}}\bigr)
  \bm C^\top,
  \qquad
  \bm C=\bigl[c_0\bm H_0^{-1},\dots,c_{K-1}\bm H_{K-1}^{-1}\bigr],
\end{equation}
where $\btV^{\mathrm{data}}$ is the covariance across observations of
the stacked conditional influence functions
$\E[\bm m_{ib}(\lambda_j)\mid Y_i,\bm W_i]$ and
$\btV^{\mathrm{sim}}$ their conditional covariance given the data.
Both matrices admit closed-form expressions under the Gaussian error
model (Gaussian tilt identities) and are estimated from the same
replicates that produce the SIMEX curve, at negligible additional
cost; a pairs-bootstrap option is provided for small samples. Two
features distinguish \eqref{eq:sandwich} from the earlier draft's
decomposition $\Var_{\mathrm{param}}+\Var_{\mathrm{extrap}}
+\Var_{\mathrm{sim}}$ with $\Var_{\mathrm{param}}
=\bm I_W(\bet_0)^{-1}$: (i)~the naive Fisher information
$\bm I_W(\bet_0)^{-1}$ does \emph{not} enter additively --- the
estimator is a fixed weighted average of the grid-point
pseudo-estimators, and its variance is the propagated influence
covariance \eqref{eq:sandwich}, which reduces to the naive variance
only in the limit of perfectly correlated grid-point influence; and
(ii)~the simulation noise enters at exactly $O(B^{-1})$, so $B=100$
makes it negligible relative to the data component --- confirmed
empirically in Section~\ref{sec:sims-val}, where the analytic sandwich
recovers $66$--$91\%$ of the parametric-bootstrap variance
(first-order under-statement, improving with $n$). Confidence
intervals are $\hat\beta_{\asimex,j}\pm
z_{1-\alpha/2}\,\widehat{\mathrm{se}}_j$ with
$\widehat{\mathrm{se}}$ from \eqref{eq:sandwich}; under
overdispersion the same sandwich is valid with robust scores, and
Section~\ref{sec:sims-nb} verifies this under negative-binomial
outcomes.

% ---------------------------------------------------------------------------
% (5) REPLACES SECTION 7.4 ("Computational Complexity") AND 7.5
%     ("Parallelization Strategy") IN FULL. The draft's eq. (35)-(36)
%     ("30-60 minutes") over-estimated by orders of magnitude because it
%     did not exploit batching or the prefix property of the CV grid.
% ---------------------------------------------------------------------------
\subsection{Computational Complexity}\label{sec:complexity}

Three implementation ideas reduce the cost far below naive estimates.
First, all $B$ SIMEX replicates at a grid point are fitted by one
\emph{batched} IRLS sweep: the score and Hessian accumulations are
tensor contractions over a $(B,n,p)$ array, so the $B$ fits share the
linear-algebra machinery. Second, the CV procedure exploits the
\emph{prefix property} of the SIMEX grid: one union-grid curve per
fold (with noise shared across $\lambda$ within a replicate, per the
manuscript's Algorithm~1) serves all candidate $\lambda_{\max}$,
reducing the number of fitted pseudo-samples per fold from
$|\mathcal K_\lambda|\cdot|\Lambda|\,B_{\mathrm{cv}}$ to
$|\Lambda|\,B_{\mathrm{cv}}$. Third, fits along the curve are warm
started from the previous grid point. Measured single-core times for
the full A-SIMEX fit (CV plus final curve plus sandwich) are $0.10$,
$0.28$, $0.95$, and $3.3$\,seconds at
$(n{=}200,p)\in\{5,15,30,50\}$ with $B=100$, $B_{\mathrm{cv}}=25$;
standard SIMEX is about twice as fast, and the corrected-score
estimator is one to two orders of magnitude cheaper but numerically
fragile (Section~\ref{sec:sims3}). The asymptotic complexity per
replication remains $O(B\,|\Lambda|\,n p^2)$ for the final curve plus
$O(5\,|\Lambda|\,B_{\mathrm{cv}}\,n p^2)$ for the CV, but the
constants are small enough that parallelization over Monte Carlo
repetitions, CV folds, or grid points---all embarrassingly
parallel, and all supported by the package through the
\texttt{future}/\texttt{furrr} backends---is a matter of convenience
rather than necessity on modern hardware.

% ---------------------------------------------------------------------------
% (6) REPLACES SECTION 8 ("Discussion") IN FULL
% ---------------------------------------------------------------------------
\section{Discussion}\label{sec:discussion}

\subsection{Summary of Contributions}

This paper makes SIMEX for Poisson regression automatic, gives it
complete and corrected asymptotic theory, and reports its
finite-sample behavior at replication standard. Three findings that we
believe are of independent interest emerged from holding the method to
this standard: (i)~held-out prediction on the observed scale cannot
select the extrapolation that bias correction requires, because its
population minimizer is the naive pseudo-parameter; (ii)~the variance
of the extrapolated functional is governed by the norm of the
extrapolation weights ($\sum_j|c_j|\approx7.8$ for the standard
quadratic extrapolant), which quantifies a substantial and previously
unappreciated efficiency premium; and (iii)~the analytic
data-plus-simulation sandwich makes inference nearly free, and its
$66$--$91\%$ recovery of the bootstrap variance delineates what part
of the finite-sample noncoverage is variance miscalibration and what
part is residual extrapolation bias.

\subsection{Comparison with Alternatives}

\emph{Versus standard SIMEX.} A-SIMEX removes the ad-hoc choices and
adds the variance machinery; in the simulations the two track each
other closely, with A-SIMEX (1-SE rule) marginally better at small
error and marginally worse at large error, where the parsimony rule
reinforces the CV's under-extrapolation. The honest summary is that
the tuning layer is not where the accuracy is won or lost; the
extrapolation bias is.

\emph{Versus the corrected score (Nakamura, 1990).} The corrected
score is consistent for the Gaussian-error pseudo-score without
requiring a model for $\Pr(X\mid W)$ --- it requires only the error
distribution, as SIMEX does --- and is dramatically cheaper
($<0.05$\,s at $p=50$). Its weaknesses, quantified in
Section~\ref{sec:sims}, are numerical: Newton divergence in about
$1\%$ of fits under misspecification, RMSE exploding with $p$ when
many coefficients are weakly identified ($0.13$ at $p=5$ to $2.72$ at
$p=50$), and instability at small $n$ with large error. SIMEX-type
methods are more stable in exactly these regimes, at a
computational premium. The two are complements: the corrected score
when $p$ is large and the error small, SIMEX when stability matters.

\emph{Versus deconvolution and functional methods.} These remain
computationally intensive and ill-posed; A-SIMEX scales to $p=50$
in seconds. A fair comparison at the level of
\citet{delaigle2008} would be valuable but is beyond our scope.

\subsection{Limitations and Caveats}

\begin{enumerate}\itemsep3pt
\item \textbf{Error variance known or estimated.} The theory
      conditions on $\bSig_\epsilon$. The replicate-based estimator
      performs indistinguishably in Section~\ref{sec:sims-sigma2};
      the estimated-variance asymptotics ( uncertainty propagation
      through the extrapolation) are not developed here.
\item \textbf{Classical error.} The NHANES application makes the
      assumption auditable: the corrected estimates lie
      $1.5$ standard errors \emph{above} the measured-weight anchor,
      the signature of nonclassical (truth-correlated) self-report
      error, which no classical-error method removes. Users with
      validation data should always fit the anchor when it exists.
\item \textbf{Residual extrapolation bias at large error.} At
      $\sigma^2_\epsilon/\Var(X)\gtrsim0.5$, polynomial SIMEX of any
      flavor retains bias of the same order as the naive attenuation
      it removes; coverage of nominal $95\%$ intervals degrades
      accordingly. Remedies (deconvolution-aware selection;
      variance-penalized degree selection) are stated in
      Section~\ref{sec:sims-val} but not developed.
\item \textbf{First-order variance estimator.} The sandwich
      under-states the finite-sample variance by $10$--$34\%$ at
      $n\le100$; bootstrap standard errors are the recommended
      small-sample alternative.
\item \textbf{High dimensions.} Repeated GLM fitting remains
      $O(np^2)$ per pseudo-sample; regularized SIMEX would be needed
      for $p$ in the hundreds.
\end{enumerate}

\subsection{Extensions and Future Work}

Beyond the Berkson error, imputation with validation data,
alternative GLMs, time-series/spatial dependence, and regularized
extensions sketched in the earlier draft (all still open), the present
results add two concrete problems: (i)~\emph{deconvolution-aware model
selection} --- constructing a held-out criterion whose population
minimizer is $\beta(-1)=\bet_0$ rather than $\beta(1)$, e.g.\ by
evaluating the held-out deviance at pseudo-deconvolved covariates
simulated under each candidate hyperparameter --- for which our
CV-bias finding provides both the motivation and the diagnostic; and
(ii)~\emph{variance-aware degree selection} --- enlarging the degree
grid while penalizing the extrapolation variance
$\dot{\bm C}_\beta$ of the efficiency theorem, in the spirit of
model-selection penalties, which would automate the bias--variance
trade-off that practitioners currently negotiate by hand.

% ---------------------------------------------------------------------------
% (7) REPLACES SECTION 9 ("Conclusion") IN FULL
% ---------------------------------------------------------------------------
\section{Conclusion}\label{sec:conclusion}

Measurement error in covariates is pervasive, and the naive Poisson
fit is not merely inefficient but misleading: its attenuation grows
linearly in the error variance and its intervals collapse, as our
simulations document cell by cell. SIMEX remains the most practical
correction framework for this problem, and this paper makes it
automatic, gives it correct and complete asymptotics, equips it with a
closed-form variance estimator, and reports---at replication
standard, with all seeds fixed---exactly where it succeeds (small to
moderate error, with valid inference and RMSE ties with the best
competitor), what it costs (a quantified efficiency premium relative
to the oracle), and where it fails (large error, where residual
extrapolation bias dominates and observed-scale cross-validation
under-extrapolates by construction). The NHANES application shows the
workflow we recommend in practice: estimate the error variance from
measurable discrepancies, correct, validate against any available
anchor, and report the sensitivity profile over the error variance.
The open-source implementation and the fully reproducible archive are
offered so that both the method and the standard of evidence can be
checked and extended.

% ===========================================================================
% (8) CHANGE LOG: superseded claims of the June 2026 draft
% ===========================================================================
% ---------------------------------------------------------------------------
% \section*{Change Log (editorial; remove before submission)}
%
% 1. Abstract / Sec 1.2(3) / Sec 8.1(3): "variance decomposes into
%    parameter estimation, extrapolation, and simulation components"
%    -> REPLACED by the exact data + O(1/B) simulation decomposition
%    (eq. sandwich); the naive Fisher information does not enter
%    additively. Evidence: corrected-proofs Theorem 2 and Remark;
%    Section 5.7 (66-91% bootstrap recovery).
%
% 2. Abstract / Sec 5.4 "expected results": "A-SIMEX bias
%    substantially smaller than standard SIMEX, approaches oracle"
%    -> NOT SUPPORTED. Real M=1000 results: at sigma^2=1.0, std SIMEX
%    bias -0.149 vs A-SIMEX(1-SE) -0.191 (n=1000); the CV criterion's
%    population minimizer is beta(1), not beta_0. Now reported as a
%    finding with remedies (Sec 5.1, 5.7). Evidence: Table 1, Fig 1,
%    CV-pick distribution (argmin picks (0.5,2) in 1000/1000 reps).
%
% 3. Sec 1.2(3) / Sec 8.1(3): "relative efficiency bounds ... ARE in
%    [0.7,0.95]" -> REPLACED. Real RMSE ratio to oracle ~2.6 at
%    (n=1000, sigma^2=0.1); the second-order efficiency comparison is
%    design-dependent in sign (corrected-proofs Theorem 3 and the
%    linear-model counterexample). Evidence: Table 1.
%
% 4. Sec 8.2 "Corrected score ... require specification of Pr(X|W)":
%    INCORRECT. The corrected score requires the error distribution
%    only, as SIMEX does. Comparison rewritten around its actual
%    weaknesses (numerical fragility; RMSE growth in p). Evidence:
%    Tables 2-3.
%
% 5. Sec 2.6 / Sec 8.2: "remains valid under mild misspecification"
%    -> now QUANTIFIED (Laplace DGP, Table 2, Fig 2); corrected-score
%    divergence 1.0% documented.
%
% 6. Sec 6 (crime data): "template; dataset to be chosen" -> REPLACED
%    by the NHANES 2017-18 application with measurable error, an
%    available gold-standard anchor, and a sigma^2 sensitivity
%    analysis (results/section6_application.tex).
%
% 7. Sec 7.4-7.5: "30-60 minutes per fit" -> REPLACED by measured
%    0.10-3.3 s (batched IRLS + prefix CV + warm starts); the draft's
%    complexity constant ignored batching. Evidence: Table 3.
%
% 8. Sec 4 (theory): proof sketches superseded in full by the
%    corrected-proofs document: the identity W(-1)=W-eps=X was wrong
%    (simulated vs observed noise); ineq. (43) (polynomial
%    approximation) was false; the B/n->0 and B/sqrt(n)->infty
%    conditions were unmotivated (B->infty suffices); Theorem 6/7/8
%    renumbered and restructured with an added growing-degree CLT and
%    an information-monotonicity lemma replacing the false
%    positive-semidefiniteness claim. Evidence: asimex_proofs.pdf,
%    its Section 8 change list.
%
% 9. Sec 3.3: Var_param + Var_extrap + Var_sim decomposition replaced
%    by eq. (sandwich); parametric bootstrap retained as a small-
%    sample option. Evidence: corrected-proofs Corollary 1;
%    validate_sandwich.py.
%
% 10. Sec 5.3 "p up to 50" timing claims and Sec 7.5 parallelization
%     strategy: revised to measured numbers; prefix property added to
%     the algorithm description (shared noise across lambda).
% ---------------------------------------------------------------------------
